% TRANSLATION-PRODUCER DRAFT — UNCHECKED
% Work: Noether Paper 42, Korean U04
% Source snapshot lines: 76--88
% Source snapshot SHA-256: B6BB3A6267BA8495FC19914A72768351E4923B13374634701AF3CBDE659883CC
% Historical whole-authority SHA-256: 443EF950D7D45DC6D9E44A9B87501620C10DA873E50E5F2B253ECCAE6A946D27
% No source check, Korean review, compilation, rendering, assembly, or approval performed.

\medskip
\noindent\textbf{3. ---} § 3을 염두에 두고, $K$가 분할되지 않는 경우에도 $K$에서 $\bar K$로 옮겨 갈 수 있음을 덧붙이고자 한다\srcfn{(6)}{이로부터 비갈루아의 경우에 대한 교차곱 표현의 유사물을 얻을 수 있다. 이는 아마도 한 학위논문에서 더 전개될 것이다.}. $k$뿐 아니라 $\bar k$도 분해체이므로 $K$는 $\bar K$, 곧 $\bar\Gg$와 $\bar k$로 이루어진 교차곱과 상사이다. 이때 $k$가 분해체이므로 $k$에 대응하는 부분군 $\Hh$에는 인자계 $1$이 대응한다. 따라서 멱등원 $E_\Hh$는 2에서와 같이 정의된다. 이를 거꾸로 적용하면 $K$가 $E_\Hh\bar K E_\Hh$와 동형임을 얻는다. 실제로 이 대수는 왼쪽 아이디얼 $\bar K E_\Hh$의 자기준동형환과 동형이므로 $\bar K$와 상사이다. 그 계수는 $K$의 계수와 일치한다. 왜냐하면 $\bar K E_\Hh$의 $\bar k$ 위 계수는 $n$이고($\bar u_{S_i}^{-1}E_\Hh$들은 $\bar k$ 위에서 독립이다), $\bar K$의 $\bar k$ 위 계수는 $nh$이기 때문이다. 이는 $\bar K$가 $\bar K E_\Hh$와 동형인 $h$개의 아이디얼로 분해된다는 뜻이며, 자기준동형환에서는 차수 $h$의 전행렬대수를 곱하는 데 해당한다.

$K$, 따라서 $\bar K$가 분할되는 경우에는 이렇게 하여 2의 결과를 다시 얻는다. 실제로
\[
\begin{aligned}
E_\Hh\bar K E_\Hh
 &=E_\Hh\bar kE_{\bar\Gg}\bar kE_\Hh
 =(E_\Hh\bar kE_\Hh)E_{\bar\Gg}(E_\Hh\bar kE_\Hh)\\
 &=\Sp(\bar k/k)E_\Hh E_{\bar\Gg}E_\Hh\Sp(\bar k/k)
 =kE_\Gg k.
\end{aligned}
\]
